\documentclass[12pt]{article}
\usepackage{graphicx}
\usepackage{fancybox}
\pagestyle{headings}

%%%%%%%%%%%%%%%%%%%%%%%%%
\pagestyle{empty}
\voffset = -8mm
\textheight = 225mm 
\hoffset = -14mm 
\textwidth = 155mm 
\evensidemargin =+1.5cm  
\oddsidemargin  =+1.5cm
\parindent8mm
\frenchspacing
%%%%%%%%%%%%%%%%%%%%%%%%%

\def\R{{\bf R}}


\begin{document}

\begin{center}
\centerline{\large{\bf FIRENZE 2009}}
\vspace*{0.1truein}
\centerline{\Large{\bf WORKSHOP on}}
\vspace*{0.05truein}
\centerline{\Large{\bf CALCULUS of VARIATIONS and PDEs}}
%\vspace*{0.1truein}
%\centerline{\small{\bf Lista abstracts}}
%\vspace*{0.1truein}
%\centerline{\large{\bf FIRENZE 2009}}
\vspace*{0.25truein}
%
%
 \begin{tabular}{c}
  \includegraphics[scale=0.7]{logo.eps} \\
  \\
  %\textbf{FIRENZE 2009}\\
  %\vrule height 0 pt depth 0.2 pt width 2 in \\
  %{\textbf{Registro delle attivit\`a didattiche}}\\
 \end{tabular}
\end{center}
 %\vskip 5em

\centerline{\small{\bf Programma}}
\vspace*{0.25truein}



%\centerline{\large{\bf FIRENZE 2009}}
%\vspace*{0.1truein}
%\centerline{\Large{\bf WORKSHOP on}}
%\vspace*{0.05truein}
%\centerline{\Large{\bf CALCULUS of VARIATIONS and PDEs}}
%\vspace*{0.1truein}
%\centerline{\small{\bf Programma}}
%%\vspace*{0.035truein}
%%\centerline{\small{\today}}
%\vspace*{0.35truein}
%
%

\begin{itemize}

\centerline{\bf Venerd\`i 23 gennaio 2009}

\item{ore 13.45: {\bf Apertura del Workshop}} 

\item{ore 14.00: {\bf Bernard Dacorogna} (EPFL Lausanne),}
\emph{``Two generalizations of Darboux theorem''}. 
%(joint work with S. Bandyopadhyay)

\item{ore 14.30: {\bf Chiara Leone} (Universit\`a di Napoli `Federico II'),}
\emph{``Partial regularity for polyconvex functionals depending on 
the Hessian determinant''}.

\item{ore 15.00: {\bf  Guido De Philippis} (Universit\`a di Firenze),}
\emph{``Una dimostrazione della minimalit\`a del cono di Simons''}.

\item{ore 15.30-16.00: {\bf Pausa Caff\`e}.}

\item{ore 16.00: {\bf Arrigo Cellina} (Universit\`a di Milano-Bicocca),} 
\emph{``A special case of constructing a barrier''}.

\item{ore 16.30: {\bf Lucio Boccardo} (Universit\`a `La Sapienza' Roma),}
\emph{``Altri sviluppi di quel che Guido Stampacchia ha insegnato a 
Paolo e a me: una teoria tipo Calderon-Zygmund   per `minimi' di 
energia infinita di alcuni funzionali integrali''}.

\item{ore 17.00: {\bf Nguyen Tien Khai} (Universit\`a di Padova),}
\emph{``Quantitative Isoperimetric and Sobolev inequalities for a class of
nonconvex sets and functions''}.

\item{ore 17.30: {\bf Agnese Di Castro} (Universit\`a `La Sapienza' Roma),}
\emph{``Elliptic problems for some anisotropic operators''}.

\item{ore 18.00-19.00: {\bf Riunione per l'organizzazione del progetto 
PRIN}}


\vspace*{0.5truein}

\newpage

\centerline{\bf Sabato 24 gennaio 2009}



\item{ore 9.30: {\bf Carlo Sbordone} (Universit\`a di Napoli `Federico II'),}
\emph{``Omeomorfismi di Sobolev e sistemi di equazioni a derivate 
parziali di tipo ellittico''.}

\item{ore 10.00: {\bf  Luisa Moschini} (Universit\`a `La Sapienza' Roma),}
\emph{``Improving $L^2$ estimates to Harnack inequalities''}.

\item{ore 10.30-11.00: {\bf Pausa Caff\`e}.}

\item{ore 11.00: {\bf Paolo Maria Mariano} (Universit\`a di Firenze),}
\emph{``Ground states of quasicrystals with magnetic order''}.

\item{ore 11.30: {\bf Menita Carozza} (Universit\`a degli Studi del Sannio),}
\emph{``Vector valued minimizers of anisotropic functionals: fractional 
differentiability and estimate for the singular set''}.


\item{ore 12.00: {\bf  Marco Mazzola} (Universit\`a di Milano Bicocca),}
\emph{``Necessary conditions for solutions to variational problems''}.


\item{ore 12.30-15.00: {\bf Pranzo}.} 


\item{ore 15.00: {\bf Alessandro Ferriero} 
(Universidad Aut\'onoma de Madrid),} 
\emph{``Perimetro minimo e convessit\`a per insiemi del piano''}.




\item{ore 15.30: {\bf  Flavia Giannetti} 
(Universit\`a di Napoli `Federico II'),}
\emph{``Regularity results for Jacobian determinants''}.

\item{ore 16.00: {\bf Anna Verde} (Universit\`a di Napoli `Federico II'),}
\emph{``Everywhere regularity of functionals with $\phi$-growth''}.

\item{ore 16.30: {\bf Giuseppe Rosario Mingione} (Universit\`a di Parma),}
\emph{``Aspetti non-lineari della teoria di Calderon-Zygmund''}.

\end{itemize}

\vfill

\centerline{\bf Comitato Organizzatore}
\noindent
Il Workshop \`e organizzato dal gruppo di ricerca 
\emph{Calculus of Variations and PDE} dell'Universit\`a di Firenze: 
Giovanni Cupini,
Matteo Focardi,
Francesco Maggi,
Paolo Marcellini,
Elvira Mascolo,
Emanuele Paolini,
Gloria Papi,
Vincenzo Vespri.

\end{document}





